\section{Level-set risk framework}
\label{sec:framework}

\subsection{An explicit working posterior}
\label{sec:posterior}

For a fixed image, write \(p=(p_1,\ldots,p_N)\). We use a shared random
threshold to construct a law over masks.

\begin{definition}[Level-set working posterior \(Q_p\)]
\label{def:levelset}
Let \(T\sim\operatorname{Unif}(0,1)\) and define
\begin{equation}
  Y_T=\{i:p_i\geq T\},
  \qquad
  Y_t=\{i:p_i\geq t\}.
  \label{eq:levelset}
\end{equation}
The distribution of \(Y_T\) is denoted by \(Q_p\).
\end{definition}

The construction satisfies
\begin{equation}
  \Prb_{Q_p}(i\in Y_T)=p_i,
  \qquad
  \Prb_{Q_p}(i,j\in Y_T)=\min(p_i,p_j).
  \label{eq:levelset-marginals}
\end{equation}
Thus \(p\) is exactly the one-pixel coverage function of \(Q_p\), while the
common threshold selects the comonotone coupling of those marginals. Samples
are nested: if \(s<t\), then \(Y_t\subseteq Y_s\). This is a classical
coverage-function random set \citep{goodman,molchanov}; related shared-level
constructions appear in uncertain-contour analysis \citep{poethkow,bolin}.

Equation~\eqref{eq:levelset-marginals} is an identity internal to \(Q_p\),
not evidence that \(p_i\) equals the true marginal or that true labels are
comonotone. The working posterior is attractive because it generates coherent,
nested candidates cheaply; it can be inaccurate when ambiguity changes topology,
moves components independently, or is poorly represented by the probability map.

\subsection{Expected loss as a threshold integral}

\begin{proposition}[Threshold-integral risk]
\label{prop:integral}
For any loss satisfying Equation~\eqref{eq:bounded-assumption},
\begin{equation}
  r_{L,Q}(x)
  :=\Ex_{Y\sim Q_{p(x)}}\!\left[L_X(Y,\Yhat(x))\right]
  =\int_0^1 L_X(Y_t,\Yhat(x))\,dt .
  \label{eq:integral}
\end{equation}
\end{proposition}
This follows immediately by integrating over the uniform threshold \(T\).
The identity is exact for the stated working posterior and does not depend on
the algebraic form of \(L\). It does not assert equality with the true
conditional risk.

Let \(H_{L,x}(t):=L_X(Y_t,\Yhat(x))\). For nodes \(t_m\in(0,1)\) and
nonnegative weights \(w_m\) summing to one, define
\begin{equation}
  r_{L,M}(x)=\sum_{m=1}^M w_m H_{L,x}(t_m),
  \qquad
  C_{L,M}(x)=-r_{L,M}(x).
  \label{eq:quad}
\end{equation}
Every candidate is compared with the fixed deployed mask. By contrast,
averaging \(L(Y_{t_m},Y_{t_n})\) over pairs of level sets estimates posterior
dispersion, not conditional risk of the deployed action.

Our default is equal-width midpoint quadrature,
\begin{equation}
  t_m=\frac{m-1/2}{M},
  \qquad
  w_m=\frac1M,
  \qquad m=1,\ldots,M.
  \label{eq:midpoint-nodes}
\end{equation}
On a finite image \(H_{L,x}\) is a bounded step function and therefore has
finite total variation \(V(H_{L,x})\). The one-dimensional Koksma--Hlawka
inequality and midpoint star discrepancy give
\begin{equation}
  \left|r_{L,M}(x)-r_{L,Q}(x)\right|
  \leq \frac{V(H_{L,x})}{2M}.
  \label{eq:koksma}
\end{equation}
We use midpoint nodes as a fixed numerical rule, not because the bound proves
they optimize finite-sample AURC \citep{niederreiter1992}.

\subsection{Numerical error and posterior discrepancy}
\label{sec:error}

Let \(P_x=\Prb(Y\in\cdot\mid X=x)\) denote the true label posterior. For the
deployed action define the loss-specific discrepancy
\begin{equation}
  \Delta_{L,\Yhat}(x)
  :=
  \left|
    \Ex_{Y\sim P_x}L_X(Y,\Yhat(x))
    -
    \Ex_{Y\sim Q_{p(x)}}L_X(Y,\Yhat(x))
  \right|.
  \label{eq:loss-discrepancy}
\end{equation}
This asks only whether two posteriors agree on the expected loss of the action
being scored. It is less restrictive than total-variation control in the
precise sense that the latter implies a bound on this one functional, while the
converse need not hold. For example, if
\(\TV(P_x,Q_{p(x)})\leq\delta_x\), with total variation
defined as the supremum over events, then
\begin{equation}
  \Delta_{L,\Yhat}(x)\leq D_x\delta_x.
  \label{eq:tv-sufficient}
\end{equation}
The converse need not hold.

\begin{proposition}[Conditional total-error bound]
\label{prop:total-bound}
For midpoint quadrature and any bounded measurable \(L\),
\begin{equation}
  \left|C_{L,M}(x)-C_L^\star(x)\right|
  \leq
  \frac{V(H_{L,x})}{2M}
  +\Delta_{L,\Yhat}(x).
  \label{eq:decomp}
\end{equation}
Under the stronger sufficient condition
\(\TV(P_x,Q_{p(x)})\leq\delta_x\), the second term can be replaced by
\(D_x\delta_x\).
\end{proposition}
The proof inserts \(r_{L,Q}(x)\) between the quadrature risk and true
conditional risk and applies the triangle inequality; Appendix~\ref{app:theory}
gives the full statement.

\begin{corollary}[Population total error]
\label{cor:population-bound}
If \(\Ex[V(H_{L,X})]<\infty\) and
\(\Ex[\Delta_{L,\Yhat}(X)]<\infty\), then
\begin{equation}
  \Ex\!\left[
    \left|C_{L,M}(X)-C_L^\star(X)\right|
  \right]
  \leq
  \frac{\Ex[V(H_{L,X})]}{2M}
  +\Ex[\Delta_{L,\Yhat}(X)].
  \label{eq:population-decomp}
\end{equation}
\end{corollary}

Increasing \(M\) controls only the first term. It does not change \(p\), the
coupling selected by \(Q_p\), or the posterior-discrepancy term. Moreover,
Equation~\eqref{eq:decomp} is a confidence-value bound, not automatically an
exact-ranking or AURC-regret bound. A positive pairwise margin is sufficient
for exact rank stability; approximate AURC control can instead use suitable
uniform or distributional control of the score errors. None of those
conditions follows from Equation~\eqref{eq:decomp} alone.
